\section{Introduction}

Discrete diffusion \citep{austin2021structured, hoogeboom2021argmax,DBLP:journals/corr/abs-2211-15089,DBLP:conf/icml/LouME24,DBLP:conf/nips/ShiHWDT24} has recently emerged as a compelling alternative to the predominant autoregressive (AR) modeling paradigm in large language models (LLMs). Unlike AR models, which generate the next token in a left-to-right fashion \citep{radford2018improving}, diffusion LLMs (dLLMs) generate text by learning to reverse a noising process that progressively corrupts a \textit{block} of token sequences. In particular, masked diffusion models (MDMs; \citealt{sahoo2024simple, ou2024your}), a subclass of discrete diffusion models that use time-dependent BERT-style \citep{devlin2019bert} masking as the forward noising process, have recently demonstrated impressive performance, with models like LLaDA \citep{llada2025} and Dream \citep{ye2025dream} matching the performance of similarly sized autoregressive LLMs. %


At generation time, MDMs begin with a fully masked sequence and iteratively unmask a fixed number of tokens at randomly sampled positions in each sampling step. As a result, they offer the potential for faster inference, as they can, in principle, generate multiple tokens in parallel using a single model call. %
Despite this promise, open-source dLLMs have, until recently, lagged behind their AR counterparts in terms of inference efficiency.
This changed with Fast-dLLM \citep{fastdllm2025}, which demonstrated that a dLLM like LLaDA \citep{llada2025} can achieve higher token throughput than similarly sized LLaMA models \citep{dubey2024llama} while maintaining competitive performance. A key component of Fast-dLLM lies in its sampling heuristic: instead of unmasking a fixed number of randomly sampled token positions, it proposes to unmask all tokens whose confidences exceed a pre-specified threshold, resulting in adaptive sampling with a variable number of unmasked tokens per step.
This success has since inspired further development of increasingly sophisticated sampling heuristics (see \Cref{sec:related} for a comprehensive overview), which has further advanced the state of the art.
These heuristics can, however, be difficult to tune, and seem to work best when imposing semi-autoregressive (semi-AR) generation, in which small blocks of tokens are unmasked sequentially \citep{arriola2025block,llada2025}.
We posit that such limitations are difficult to avoid with heuristics alone, since they are essentially handcrafted solutions to a sequential decision making problem: 
\emph{What is the optimal order in which to unmask the sequence of tokens, in order to strike a good balance between efficiency and correctness?}


We propose to investigate this question by moving beyond heuristics to stress instead a \textit{learning} based approach, which relies on a transformer-based unmasking policy trained using reinforcement learning (RL).
This approach is motivated by the above observation that unmasking in dLLMs can be viewed as a (Markovian) sequential decision making problem, and follows recent successes in reinforcement learning for post-training of language models.
However, in contrast to recent works that use RL to improve the reasoning abilities of dLLMs (e.g., \citealt{zhao2025d1}), our goal is not to improve the capacity of the underlying model, but rather to use RL to automate the discovery of adaptive sampling strategies.
We therefore treat the underlying dLLM as the \emph{environment} in which to act, not the policy, leaving it unchanged and instead focusing a policy that is parameterized as a lightweight stand-alone network.
Empirically, we demonstrate that our learned sampling policies match the performance of state-of-the-art heuristic samplers \citep{fastdllm2025} in standard generation settings, while also surpassing heuristics in certain scenarios, notably when moving away from semi-AR generation. %
In summary, we make the following contributions:
\begin{itemize}
\item We formalize sampling in dLLMs as a Markov decision process (MDP) (\Cref{sec:methods_mdp}), propose a lightweight design for the unmasking policy %
(\Cref{sec:methods_policy}), and outline our RL training pipeline based on group relative policy optimization (GRPO; \citealt{shao2024deepseekmath}) (\Cref{sec:methods_grpo}).
\item We conduct extensive experiments showing that, under our proposed RL framework, learned policies can match the performance of heuristic samplers such as Fast-dLLM (\Cref{sec:exp-bl32}), while also addressing some of the challenges heuristic samplers face outside the semi-AR regime (\Cref{sec:exp-bl256}).
\item We study the transferability of our learned samplers across different models and data domains (\Cref{sec:exp-transfer}), and provide detailed insights into stabilizing RL training and policy design (\Cref{sec:exp-ablate}).
\end{itemize}

